\documentclass[12pt]{article}
\usepackage[utf8]{inputenc}
\usepackage[T1]{fontenc}
\usepackage[margin=1in]{geometry}
\usepackage{setspace}
\usepackage{xcolor}
\usepackage[colorlinks=true, urlcolor=blue!60!black, linkcolor=blue!60!black]{hyperref}
\setlength{\parindent}{2em}
\setlength{\parskip}{1.25em}
\renewcommand{\baselinestretch}{1.0}
\begin{document}
\begin{center}
{\LARGE \textit{Bradwell v. Illinois}, 83 U.S. 130 (1873)\par}
\vspace{0.8em}
{\large Case Brief --- Four-Page Edition\par}
\end{center}
\vspace{1em}

\section{Memory Jogger}

One day after gutting the Fourteenth Amendment's Privileges or Immunities Clause in the \textit{Slaughter-House Cases}, the Supreme Court held 8--1 that Illinois could bar Myra Bradwell from practicing law because admission to a state bar is not a privilege of national citizenship, while Justice Bradley's notorious concurrence invoked ``the law of the Creator'' to confine women to the domestic sphere (\href{https://supreme.justia.com/cases/federal/us/83/130/}{83 U.S. 130}).

\section{Detailed Case Facts}

Myra Bradwell, a longtime Chicago resident and the founder and editor of the influential \textit{Chicago Legal News}, applied to the Illinois Supreme Court for a law license, with a certificate of good moral character and of her legal qualifications. Her affidavit stated that she had been born in Vermont, was a U.S. citizen, and had lived in Chicago for many years. She claimed a constitutional entitlement to a license under Article IV, \S{} 2 and the Fourteenth Amendment's Privileges or Immunities Clause.

The Illinois statute required a license from two justices of the state supreme court, with a county-court certificate of good moral character as its only express prerequisite. The court first refused Bradwell on the ground of coverture: as a married woman she could not be bound by the contracts between attorney and client. After her printed argument, it held that because Illinois had adopted the English common law, under which female attorneys were unknown, the legislature never contemplated admitting women. The Court decided \textit{Bradwell} on April 15, 1873, one day after the \textit{Slaughter-House Cases}.

\section{Procedural History}

The Illinois Supreme Court denied Bradwell's application. She sued out a writ of error to the U.S. Supreme Court under \S{} 25 of the Judiciary Act of 1789, as amended, claiming violations of Article IV, \S{} 2 and the Fourteenth Amendment. The case was argued January 18, 1873, and the Court affirmed, 8--1.

\section{Judicial Votes}

Affirmed, \textbf{8--1}.
\begin{itemize}
\item \textbf{Majority (5):} Justice Miller, joined by Justices Clifford, Davis, Strong, and Hunt.
\item \textbf{Concurrence in the judgment (3):} Justice Bradley, joined by Justices Swayne and Field.
\item \textbf{Dissent:} Chief Justice Chase, without opinion; he was gravely ill and died weeks later.
\end{itemize}

\section{Holding}

Article IV, \S{} 2 protects only citizens of other states and offers no protection against the laws of one's own state; Bradwell's Chicago residence made her a citizen of Illinois. The right to admission to practice in state courts is not a privilege or immunity of \textit{United States} citizenship under the Fourteenth Amendment, so states retain full authority to set bar qualifications. Miller's remark about admission to federal courts is dicta, and Bradley's declarations about women appeared only in a three-Justice concurrence, never in the Court's holding.

\section{Analysis of Opinions}

\textbf{Majority (Miller).} In a brief, strictly doctrinal opinion that avoided gender stereotypes, Miller rejected the Article IV claim because the Fourteenth Amendment's first sentence makes citizens citizens of the state where they reside. He agreed that some privileges belong to U.S. citizens as such but denied that state bar admission is one: it had never depended on U.S. citizenship, and noncitizens had been admitted to state and federal courts. Invoking \textit{Slaughter-House} as conclusive, he found the regulation of law licenses a power never transferred to the federal government.

\textbf{Bradley (concurring in the judgment).} Bradley, Swayne, and Field had dissented in \textit{Slaughter-House}, insisting that the Clause protects the right to pursue a lawful occupation. To explain why New Orleans butchers had that right but a qualified female lawyer did not, Bradley asserted that the right had never been a fundamental privilege of women. Invoking civil law, common law, ``nature herself,'' and coverture, he dismissed unmarried women as ``exceptions to the general rule'' and held that excluding women from offices requiring ``decision and firmness'' fell within the police power.

\textbf{Chase (dissenting).} Too ill to write, Chase dissented ``from the judgment of the Court and from all the opinions,'' consistent with his support for abolition and women's suffrage.

\section{Examples: Future Applications}

\begin{enumerate}
\item \textbf{Same-Side: Hair-Braiding License.} A resident challenges a 1,500-hour cosmetology requirement for braiders solely under the Privileges or Immunities Clause. Under \textit{Bradwell} and \textit{Slaughter-House}, the claim fails: a state-regulated occupation is not a privilege of national citizenship.
\item \textbf{Same-Side: Vermont-Born Texan.} A lifelong Texas resident born in Vermont challenges a Texas medical-licensing rule under Article IV, \S{} 2. The clause is inapplicable because she is a citizen of the state whose law she challenges.
\item \textbf{Opposite-Side: Sex-Based Exclusion under Equal Protection.} A state nursing school excludes qualified applicants by sex. Modern intermediate scrutiny invalidates classifications grounded in ``archaic and stereotypic notions,'' as \textit{Mississippi University for Women v. Hogan} (1982) held.
\item \textbf{Opposite-Side: Nonresident Lawyer.} A state bars all nonresidents from its bar. An out-of-state attorney has precisely the posture Bradwell lacked, and \textit{Supreme Court of New Hampshire v. Piper} (1985) protects law practice under Article IV, \S{} 2.
\item \textbf{Fence-Sitter: Federal Bar Admission.} A federal district court excludes a class of qualified citizens on a non-suspect basis. Miller suggested federal-court admission ``would relate to citizenship of the United States,'' yet noted that noncitizens had been admitted.
\end{enumerate}

\section{Critique}

\textbf{Progressive and feminist views.} \textit{Bradwell} emblemizes how nineteenth-century courts constitutionalized patriarchal ideology. Miller's refusal to police state occupational exclusions left women without federal recourse, and Bradley raised ``separate spheres'' and coverture into ``the law of the Creator,'' ignoring Bradwell's own career. Because no Equal Protection argument was pressed, meaningful scrutiny of sex discrimination waited until \textit{Reed v. Reed} (1971).

\textbf{Originalist and textualist views.} Timothy Sandefur and Justice Thomas in \textit{McDonald} argue that \textit{Slaughter-House} and \textit{Bradwell} mutilated the Clause's original meaning. Originalists divide over the outcome: traditionalists note that ratifiers assumed sex-based exclusions permissible, while Steven Calabresi and Julia Rickert argue the Amendment barred caste legislation. Textualists fault Bradley for abandoning text for natural-law assumptions.

\section{Key Quotations}

\begin{quote}
``Man is, or should be, woman's protector and defender. The natural and proper timidity and delicacy which belongs to the female sex evidently unfits it for many of the occupations of civil life. The Constitution of the family organization, which is founded in the divine ordinance as well as in the nature of things, indicates the domestic sphere as that which properly belongs to the domain and functions of womanhood.''
\end{quote}
\noindent --- Justice Bradley, concurring in the judgment, 83 U.S. at 141.

\begin{quote}
``The paramount destiny and mission of woman are to fulfil the noble and benign offices of wife and mother. This is the law of the Creator. And the rules of civil society must be adapted to the general constitution of things, and cannot be based upon exceptional cases.''
\end{quote}
\noindent --- Justice Bradley, concurring in the judgment, 83 U.S. at 141--42.

\begin{quote}
``We agree with him that there are privileges and immunities belonging to citizens of the United States, in that relation and character, and that it is these and these alone which a state is forbidden to abridge. But the right to admission to practice in the courts of a state is not one of them. This right in no sense depends on citizenship of the United States.''
\end{quote}
\noindent --- Justice Miller, for the Court, 83 U.S. at 139.

\end{document}
